\documentclass[a4paper,12pt]{article}
\usepackage[T2A]{fontenc}
\usepackage[utf8]{inputenc}
\usepackage[russian]{babel}
\usepackage{geometry}
\usepackage{amsmath, amssymb}
\geometry{margin=2.5cm}

\title{Task 2.2}

\begin{document}

\section*{Замечания}
Через $A$ $<$ $B$ обозначаем что событие $A$ произошло раньше события $B$.

\begin{itemize}
\item Между $A.1$ и $A.5$ A удерживает $A\_flag = true$.
\item $B\_flag$ изменяется в $B.1$, $B.3$, $B.5$, $B.7$.
\item $A\_flag$ изменяется в $A.1$ и $A.5$.
\end{itemize}

\section*{2.2.a}

\textbf{Доказательство} что оба потока не могут одновременно находиться в critical\_section().

\subsection*{Доказательство от противного}

\begin{itemize}

\item Пусть существует такая трасса, в которой потоки одновременно находятся в критической секции.

\item Тогда существуют события входа в секцию: $A.4$ и $B.6$. И они наступают в одно и то же время.

\item Без потери общности предполагаем $A.4$ < $B.6$.

\item Чтобы поток A выполнил $A.4$, он должен выйти из цикла \texttt{while (is\_raised\_B())}.
Значит, существует последнее чтение
$A.2^{\ast}$ перед $A.4$,
которое вернуло \texttt{false}.
Значит в момент $A.2^{\ast}$: B\_flag = false.


\item Чтобы B выполнил $B.6$, он должен выйти из цикла \texttt{while (is\_raised\_A())}.
Значит, существует последнее чтение
$B.2^{\ast}$ перед $B.6$,
которое вернуло \texttt{false}.
Значит в момент $B.2^{\ast}$: A\_flag = false.

\item Между $A.1$ и $A.5$ $A\_flag = true$, а $A.5$ выполняется после выхода из критической секции.

\item Поскольку интервалы перекрываются
и $A.4$ < $B.6$,
имеем: $B.6$ < $A.5$.

\item Тогда $B.2^{\ast}$ < $B.6$ < $A.5$.

\item Следовательно, в момент $B.2^{\ast}$ $A\_flag$ обязан быть равен \texttt{true}.
А это противоречит тому, что чтение вернуло \texttt{false}.

\item Получено противоречие.

\end{itemize}

Следовательно, алгоритм обеспечивает взаимное исключение.

\section*{2.2.b}

\textbf{Доказать:} что невозможно, чтобы оба потока навсегда зависли, не входя в критическую секцию.

\subsection*{Доказательство от противного}

\begin{itemize}

\item Пусть, существует такая трасса, в которой ни $A.4$, ни $B.6$ больше не происходят и она бесконечная.

\item A может не входить в критическую секцию только оставаясь бесконечно в цикле $A.2$.
Значит он бесконечно читает $B\_flag$ и каждый раз видит \texttt{true}.

\item B может не входить в критическую секцию только оставаясь бесконечно в цикле $B.2$.

\item Если поток B видит $A\_flag = true$, он обязан выполнить $B.3$ и записать $B\_flag = false$.

\item Поскольку B бесконечно выполняется, в трассе существует событие $B.3$.

\item Поскольку трасса упорядочена, то существует следующее чтение
$A.2_{\text{next}}$ такое, что $B.3$ < $A.2_{\text{next}}$.

\item В момент $A.2_{\text{next}}$  A увидит $B\_flag = false$ и выйдет из цикла ожидания.

\item Следовательно, A выполнит $A.4$, что противоречит тому предположению что, ни один поток больше не входит
в критическую секцию.

\end{itemize}

Следовательно, алгоритм deadlock-free.

\end{document}